% Generated draft: requires mathematical and source review.
\documentclass[11pt]{article}
\usepackage[margin=25mm]{geometry}
\usepackage{fontspec}
\setmainfont{Latin Modern Roman}
\usepackage{amsmath,amssymb,mathtools,unicode-math}
\setmathfont{Latin Modern Math}
\usepackage{enumitem,needspace,xurl,hyperref}
\setlength{\parindent}{0pt}
\setlength{\parskip}{5pt}
\setlength{\emergencystretch}{3em}
\newcommand{\N}{\mathbb{N}}
\newcommand{\Z}{\mathbb{Z}}
\newcommand{\Q}{\mathbb{Q}}
\newcommand{\R}{\mathbb{R}}
\newcommand{\C}{\mathbb{C}}
\newcommand{\F}{\mathbb{F}}
\newcommand{\A}{\mathbb{A}}
\newcommand{\PP}{\mathbb{P}}
\newcommand{\E}{\mathbb{E}}
\newcommand{\eps}{\varepsilon}
\newcommand{\dd}{\,\mathrm{d}}
\begin{document}
\section*{UnsolvedMath ID: 30005398}
\textbf{Source problem code:} catalogue entry
\subsection{Definitions and mathematical statement}
\textbf{Source setting and unresolved hypotheses.} The supplied record attributes the following three conjectures to Fargues. Write \(X^{\mathrm{alg}}=X_{E,F}^{\mathrm{alg}}\) and \(X^{\mathrm{ad}}=X_{E,F}^{\mathrm{ad}}\) for the algebraic and adic Fargues--Fontaine curves named in the source. The record does not specify the hypotheses on \(E\) and \(F\), their characteristics, or the prime \(p\) appearing in its literature assessment. These parameters retain their intended source meaning; the assertions below do not quantify over arbitrary fields \(E,F\). Recovering the missing hypotheses is necessary for a complete mathematical formulation.

An étale abelian torsion sheaf \(\mathcal F\) on \(X^{\mathrm{alg}}\) is a sheaf of abelian groups on its étale site whose geometric stalks are torsion groups. Write \(H^i_{\mathrm{\acute et}}(X,\mathcal G)\) for étale sheaf cohomology. The source supplies a morphism of locally ringed spaces
\[
u:X^{\mathrm{ad}}\longrightarrow X^{\mathrm{alg}}
\]
and asserts that it induces a morphism of étale sites in the same direction. Let \(\mathcal F^{\mathrm{ad}}=u^*\mathcal F\), where \(u^*\) denotes the inverse-image functor on abelian sheaves associated with that morphism of sites. Let \(K=E(X^{\mathrm{alg}})\) denote the function field named in the source. A field \(K\) is \(C_1\) if, for every pair of integers \(d\geq 1\) and \(n>d\), every homogeneous polynomial \(f\in K[T_1,\ldots,T_n]\) of degree \(d\) has a zero in \(K^n\setminus\{0\}\).

\textbf{Requested assertions, subject to the missing source hypotheses.}
\begin{enumerate}
\item For every étale abelian torsion sheaf \(\mathcal F\) on \(X^{\mathrm{alg}}\) and every integer \(i\geq 3\),
\[
H^i_{\mathrm{\acute et}}(X^{\mathrm{alg}},\mathcal F)=0.
\]
\item The morphism \(u\) induces the stated morphism of étale sites, and for every such \(\mathcal F\) and every integer \(i\geq 0\), the resulting comparison map is an isomorphism:
\[
H^i_{\mathrm{\acute et}}(X^{\mathrm{alg}},\mathcal F)
\xrightarrow{\ \sim\ }
H^i_{\mathrm{\acute et}}(X^{\mathrm{ad}},u^*\mathcal F).
\]
\item The field \(K=E(X^{\mathrm{alg}})\) is \(C_1\).
\end{enumerate}
The first two assertions as recorded concern all étale abelian torsion sheaves, without a stated prime-to-\(p\) restriction or a requirement that one integer annihilate the whole sheaf.

\textbf{Qualifications retained from the report.} The quoted report remarks that a \(C_1\) field has cohomological dimension at most \(1\), that absolute purity and the \emph{méthode de la trace} give the implication \((3)\Rightarrow(1)\), and that results discussed in the talk give \((1)\Rightarrow(2)\). These are retained as source remarks in the intended setting; the omitted parameter hypotheses and talk context have not been reconstructed here.
\subsection{Short English statement}
Fargues' recorded conjectures ask for étale cohomology to vanish in degrees at least \(3\) for every torsion abelian sheaf on the algebraic Fargues--Fontaine curve, for comparison with the adic curve in every nonnegative degree, and for the function field to be \(C_1\). The supplied record omits the hypotheses on \(E,F\). Its dated curation labels the problem partially solved; this draft retains the full recorded targets and requires human mathematical review.
\subsection{Sources}
\begin{itemize}
\item Supplied local UnsolvedMath registry record, permanent ID \(30005398\), title ``Fargues' Cohomological Conjectures for the Curve,'' exact source problem code \texttt{catalogue entry}, in the Oberwolfach Reports Open Problems collection. The quoted conjecture is attributed to Fargues; the separate proposer field is empty, and the supplied proposed year is \(2023\).
\item Supplied report citation: \emph{Arithmetic of Shimura Varieties} (2023), Oberwolfach Reports, pp.\ 261--326, DOI \texttt{10.4171/owr/2023/5}. \url{https://doi.org/10.4171/owr/2023/5}.
\item Supplied literature source: Bartling, ``etale cohomology of the Fargues--Fontaine curve.'' \url{https://arxiv.org/abs/2206.14253}.
\item Supplied literature source: ``original Fargues conjectures in OWR.'' \url{https://ems.press/content/serial-article-files/46998?nt=1}.
\item Literature assessment dated exactly \texttt{2026-08-22}: the supplied curation labels the problem partially solved and reports that Bartling proves the prime-to-\(p\) vanishing and algebraic--adic comparison conjectures, while the \(p\)-torsion statement remains conditional and the \(C_1\) function-field assertion remains unresolved. This is supplied dated curation, not an independently verified assessment; its time-sensitive status and cited claims require independent checking.
\item The registry labels the statement \texttt{corrected\_verified} and states that it was reconstructed and checked against the official Oberwolfach report. This is a supplied verification claim, not verification performed for this draft. The excerpt also says ``also in [3],'' but the corresponding bibliography entry is absent from the supplied record.
\end{itemize}
\end{document}
